% from ega4-8-fr.tex
\oldpage[IV-3]{11}
Soit $T$ l'espace topologique limite du système
$(S_\lambda,\psi_{\lambda\mu})$ et posons
$\psi=\varprojlim\psi_\lambda:S\to T$. On peut se borner au cas où
$S_0=S_\alpha$ pour un $\alpha\in L$, et $\lambda\geq\alpha$.

\medskip
\noindent
(i) Montrons en premier lieu que la topologie de $S$ est image réciproque par
$\psi$ de la topologie de $T$~; autrement dit, si
$\pi_\lambda:T\to S_\lambda$ est l'application canonique, il faut montrer
que tout ouvert de $S$ est réunion d'ouverts de la forme
$\psi^{-1}(\pi_\lambda^{-1}(U_\lambda))$, où $U_\lambda$ est ouvert dans
$S_\lambda$. Or tout ouvert de $S$ est par définition réunion d'ouverts
obtenus de la façon suivante~: on considère un ouvert affine $U_0$ de $S_0$,
d'anneau $A_0$, de sorte que $\psi_\alpha^{-1}(U_0)$ est l'ouvert affine de
$S$ d'anneau $A=\Gamma(U_0,\mathscr A)$, puis on prend un élément $f\in A$ et
on considère dans $\operatorname{Spec}(A)$, identifié à
$\psi_\alpha^{-1}(U_0)$, l'ouvert $D(f)$. Ce sont ces ouverts $D(f)$ qui
forment une base de la topologie de $S$ (II, 1.3.1). Or, si on pose
$A_\lambda=\Gamma(U_0,\mathscr A_\lambda)$, on a
$A=\varinjlim A_\lambda$ (I, 1.3.9), donc il existe un indice $\lambda$ tel
que $f$ soit l'image canonique d'un élément $f_\lambda\in A_\lambda$~; on a
alors $D(f)=\psi_\lambda^{-1}(D(f_\lambda))$ (I, 1.2.2), et comme
$\psi_\lambda=\pi_\lambda\circ\psi$, notre assertion est démontrée.

\medskip
\noindent
(ii) Prouvons maintenant que $\psi$ est \emph{bijective}, ce qui achèvera la
démonstration. Comme $S$ est un espace de Kolmogoroff, il résulte déjà de (i)
que $\psi$ est injective, et il reste donc à montrer que $\psi$ est
surjective. On peut évidemment remplacer pour cela $T$ par un ouvert
$\pi_\alpha^{-1}(U_0)$, où $U_0$ est ouvert affine dans $S_\alpha=S_0$, donc
on peut se limiter au cas où les $S_\lambda$ et $S$ sont affines, autrement
dit $\mathscr A_\lambda$ est le faisceau associé à une $A_0$-algèbre
$A_\lambda$, et $\mathscr A$ le faisceau d'algèbres associé à
$A=\varinjlim A_\lambda$~; nous noterons encore
$\varphi_{\mu\lambda}:A_\lambda\to A_\mu$ et
$\varphi_\lambda:A_\lambda\to A$ les homomorphismes canoniques. Par
définition, un élément de $T$ est une famille
$(\mathfrak p_\lambda)_{\lambda\in L}$, où $\mathfrak p_\lambda$ est un
idéal premier de $A_\lambda$ et où l'on a
$\mathfrak p_\lambda=\varphi_{\mu\lambda}^{-1}(\mathfrak p_\mu)$ pour
$\lambda\leq\mu$. On sait alors (5.13.3 et 5.13.1) qu'il y a un idéal premier
$\mathfrak p$ de $A$ tel que
$\mathfrak p_\lambda=\varphi_\lambda^{-1}(\mathfrak p)$ pour tout
$\lambda\in L$, ce qui achève la démonstration.

En particulier, on a ainsi démontré le

\begin{corollary}[8.2.10]
\label{IV.8.2.10-fr}
Soit $(A_\lambda)_{\lambda\in L}$ un système inductif filtrant d'anneaux, et
soient $A=\varinjlim A_\lambda$, $\varphi_\lambda:A_\lambda\to A$ les
homomorphismes canoniques. L'application canonique
$\mathfrak p\mapsto(\varphi_\lambda^{-1}(\mathfrak p))$ est un
homéomorphisme de $\operatorname{Spec}(A)$ sur l'espace topologique
$\varprojlim\operatorname{Spec}(A_\lambda)$.
\end{corollary}

\begin{corollary}[8.2.11]
\EGAReaderTarget{ega4.c04.cor.8-2-11}{IV.8.2.11-fr}
Avec les notations de (8.2.8), pour tout ouvert quasi-compact $U$ de $S$, il
existe un indice $\lambda$ et un ouvert quasi-compact $U_\lambda$ de
$S_\lambda$ tel que $U=\psi_\lambda^{-1}(U_\lambda)$.
\end{corollary}

Cela résulte de ce que, par définition de la topologie limite projective, les
$\psi_\lambda^{-1}(U_\lambda)$ ($U_\lambda$ ouvert quasi-compact de
$S_\lambda$) forment une base de la topologie de $S$, et de ce que l'ensemble
d'indices $L$ est filtrant.

\begin{corollary}[8.2.12]
\label{IV.8.2.12-fr}
Avec les notations de (8.2.8), la limite inductive du système inductif
d'homomorphismes
$\theta_\lambda^\sharp:\psi_\lambda^*(\mathcal O_{S_\lambda})\to
\mathcal O_S$ de faisceaux d'anneaux sur $S$ est un isomorphisme
\begin{equation}
\varinjlim_\lambda\psi_\lambda^*(\mathcal O_{S_\lambda})
\xrightarrow{\sim}\mathcal O_S.
\tag{8.2.12.1}\label{IV.8.2.12.1-fr}
\end{equation}
\end{corollary}

On peut évidemment supposer les $S_\lambda$ affines~; avec les notations de
la démonstration de (8.2.9), tout revient à voir que la limite inductive du
système d'applications canoniques
$(A_\lambda)_{\mathfrak p_\lambda}\to A_{\mathfrak p}$ est un isomorphisme,
ce qui n'est autre que (5.13.3, (ii)).

